% Answers to Chapter 12, from lectures/L12/appendix/c_solutions.md; the self-test from
% lectures/L12/README.md.

\facitavsnitt{c12}{Kapitel 12}{Derivata, del I}
\begin{facit}

\svar{1.1} $G_{\text{lin}} = 10^{1{,}3} \approx 20$

\svar{1.2} $P = 10^{1{,}7}\,\text{mW} \approx 50\,\text{mW}$

\svar{1.3} $n = 10^{0{,}48} \approx 3$

\svar{2.1} \sv{a} $f'(x) = -6x + 5$ \sv{b} $f'(x) = 6x^2 - 10x + \frac{2}{3}$
\sv{c} $f'(x) = 8x^3 - 3x^2 + \frac{6x}{5} + 4$

\svar{2.2} \sv{a} $f'(x) = -2x + 8$ \sv{b} $x = 4$ \sv{c} Maximum, eftersom
$f''(4) = -2 < 0$. \sv{d} Största värdet är $f(4) = 9$; maximipunkten är $(4, 9)$.
\sv{e} En nedåtvänd parabel med toppen i $(4, 9)$ och nollställena $x = 1$ och $x = 7$.

\svar{2.3} \sv{a} $i'(t) = -0{,}6t^2 + 3{,}6t - 3{,}6$
\sv{b} $t = 3 \pm \sqrt{3}$, alltså $t_1 \approx 1{,}27\,\text{s}$ och
$t_2 \approx 4{,}73\,\text{s}$.
\sv{c} $i''(t_1) \approx 2{,}08 > 0$: minimum vid $t_1$; $i''(t_2) \approx -2{,}08 < 0$:
maximum vid $t_2$.
\sv{d} Minimivärdet är $i(t_1) \approx 0{,}42\,\text{A}$ och maximivärdet
$i(t_2) \approx 4{,}58\,\text{A}$.
\sv{e} Kurvan startar i $i(0) = 2{,}5\,\text{A}$, har en dal vid $(1{,}27;\, 0{,}42)$ och en topp
vid $(4{,}73;\, 4{,}58)$ och faller sedan.

\svar{3.1} \sv{a} $5x^4$ \sv{b} $12x^2$ \sv{c} $0$ \sv{d} $4x - 5$ \sv{e} $-12x^3 + 1$
\sv{f} $x^2$

\svar{3.2} \sv{a} Kvoten förenklas till $4 + h$. \sv{b} $4$
\sv{c} $f'(2) = 2 \times 2 = 4$, samma svar. \sv{d} $4{,}1$ respektive $4{,}01$

\svar{3.3} \sv{a} $f'(x) = 3x^2 - 3$ \sv{b} $-3$ \sv{c} $9$
\sv{d} Där $x = -1$ och $x = 1$, det vill säga i $(-1, 2)$ och $(1, -2)$.

\svar{3.4} \sv{a} $2$ \sv{b} $2$ \sv{c} $y = 2x - 4$ \sv{d} I punkten $(2, 0)$.

\svar{3.5} \sv{a} $f'(x) = 2x - 6$ \sv{b} $x = 3$ \sv{c} Minimum, eftersom $f''(x) = 2 > 0$.
\sv{d} $f(3) = -1$; minimipunkten är $(3, -1)$.

\svar{3.6} \sv{a} $f'(x) = 3x^2 - 6x - 9 = 3(x - 3)(x + 1)$ \sv{b} $x = -1$ och $x = 3$
\sv{c} Maximum i $x = -1$ ($f''(-1) = -12 < 0$), minimum i $x = 3$ ($f''(3) = 12 > 0$).
\sv{d} $f(-1) = 10$ och $f(3) = -22$

\svar{3.7} \sv{a} $f'(x) = 2x - 4$ \sv{b} $x > 2$ \sv{c} $x < 2$
\sv{d} $f'(2) = 0$: funktionen går från avtagande till växande, så $x = 2$ är en minimipunkt.

\svar{3.8} \sv{a} $f'(x) = 4x^3 - 4x$, $f''(x) = 12x^2 - 4$
\sv{b} $f'(x) = 15x^2 + 2$, $f''(x) = 30x$ \sv{c} $f'(x) = -2x + 7$, $f''(x) = -2$
\sv{d} $f''(x) > 0$: kurvan är konkav uppåt, som en dal; $f''(x) < 0$: konkav nedåt, som en
kulle.

\svar{3.9} \sv{a} $u'(t) = 8t + 2$ \sv{b} $i_C(t) = (0{,}8t + 0{,}2)\,\text{mA}$
\sv{c} $1{,}0\,\text{mA}$ \sv{d} $t = 2{,}25\,\text{s}$

\svar{3.10} \sv{a} $i'(t) = 4t - 4$ \sv{b} $u_L(t) = (2t - 2)\,\text{V}$ \sv{c} $t = 1\,\text{s}$
\sv{d} Strömmen är stationär, ett minimum på $i(1) = 1\,\text{A}$: den slutar för ett ögonblick
att ändras.

\svar{3.11} \sv{a} $P'(I) = 24 - 6I$ \sv{b} $I = 4\,\text{A}$
\sv{c} Maximum, eftersom $P''(I) = -6 < 0$. \sv{d} $P(4) = 48\,\text{W}$

\svar{3.12} \sv{a} $u'(t) = -4t + 12$ \sv{b} Vid $t = 3\,\text{s}$, då $u = 18\,\text{V}$.
\sv{c} Den ökar med $8\,\text{V/s}$. \sv{d} $t = 4\,\text{s}$

\svar{3.13} \sv{a} Växande, eftersom $f'(1) > 0$. \sv{b} Avtagande, eftersom $f'(3) < 0$.
\sv{c} En maximipunkt: $f'(2) = 0$ och $f''(2) < 0$.
\sv{d} Grafen stiger fram till $x = 2$, har en vågrät tangent i toppen och faller sedan.

\svar{3.14} \sv{a} Inget extra $x$ ska stå kvar: $f'(x) = 3x^2$.
\sv{b} En konstant har derivatan noll: $f'(x) = 0$.
\sv{c} Derivatan av $4x$ är lutningen: $f'(x) = 4$.
\sv{d} Termen $2x$ ska också deriveras: $f'(x) = 6x + 2$.
\sv{e} Falskt: en stationär punkt kan också vara ett minimum eller en terrasspunkt.
\sv{f} Falskt: $f''(x)$ fås genom att derivera $f'(x)$ en gång till. Till exempel ger $f(x) = x^3$
att $f''(x) = 6x$, inte $9x^4$.

\svar[Testa dig själv]{T} \sv{1} $f'(x) = 12x^2 - 4x$
\sv{2} $(1, -9)$, ett minimum eftersom $f''(x) = 2 > 0$.

\end{facit}
